\chapter{Tests and validation}
\label{cap:testes}

\section{Strategy}

The project does not have automated unit tests in the usual sense
(\code{cargo test}), because the code depends on hardware, on a bootloader
and on a bare metal execution environment. Validation was performed through
\emph{integration tests} in QEMU, with verification of:

\begin{itemize}
  \item textual output on the serial port;
  \item contents of VGA memory, read through the QEMU monitor;
  \item input from the PS/2 keyboard, simulated with \code{sendkey};
  \item input from the serial port, sent over a Unix socket;
  \item validity of the ELF and of the ISO with external tools.
\end{itemize}

\section{Test matrix}

\begin{table}[H]
\centering
\caption{Matrix of executed tests.}
\label{tab:testmatrix}
\begin{tabularx}{\linewidth}{@{}lXc@{}}
\toprule
\textbf{ID} & \textbf{Test} & \textbf{Result} \\
\midrule
T01 & \code{make clean \&\& make} completes without errors & Passed \\
T02 & \code{grub-file --is-x86-multiboot2} accepts the ELF & Passed \\
T03 & ISO recognized as El Torito BIOS & Passed \\
T04 & QEMU/SeaBIOS initializes the ISO and reaches the banner & Passed \\
T05 & \code{rinux>} prompt appears on the VGA screen & Passed \\
T06 & PS/2 input: \code{help} lists the commands & Passed \\
T07 & PS/2 input: \code{echo ola} prints \code{ola} & Passed \\
T08 & Serial input: \code{echo serial ok} works & Passed \\
T09 & Prompt returns after each command & Passed \\
T10 & \code{info} shows the kernel configurations & Passed \\
T11 & \code{clear} sends \code{ESC[2J ESC[H} and blanks the screen & Passed \\
T12 & \code{reboot} and \code{halt} end the flow & Verified manually \\
T13 & IDT installed and timer IRQ0 increments \code{ticks} & Passed \\
T14 & \code{demand} raises \code{\#PF} and maps the kernel page & Passed \\
T15 & \code{usertest} runs in ring 3, raises two page faults and returns & Passed \\
T16 & \code{userprot} violates U/S and is terminated by the handler & Passed \\
T17 & \code{mem} shows U/S = 1 for user and 0 for kernel & Passed \\
T18 & \code{tty} reports \code{tty0} and \code{tty 3} switches console & Passed \\
T19 & \code{getty} logs in with \code{root}/\code{rinux} in ring 3 & Passed \\
T20 & Password is not echoed; a wrong password exits with status 1 & Passed \\
T21 & \code{spawn ... signal} + \code{kill <pid> USR1} runs the handler & Passed \\
T22 & Default \code{SIGTERM}/\code{SIGKILL} actions terminate the target & Passed \\
T23 & \code{SIGSTOP} shows \code{parada} in \code{ps}; \code{SIGCONT} resumes & Passed \\
T24 & A blocked signal stays pending and is delivered after unblocking & Passed \\
T25 & \code{wait} reaps a child with status 42 and reports signal status 143 & Passed \\
T26 & \code{/proc/signals}, \code{sig} and \code{sigaction} show the masks & Passed \\
T27 & \code{kill(pid, 0)} succeeds for a zombie and fails after \code{wait} & Passed \\
\bottomrule
\end{tabularx}
\end{table}

\section{Execution evidence}

\subsection{Banner and prompt}

\begin{lstlisting}[style=shell,caption={VGA screen after boot.},label={lst:testboot}]
==============================================================
  Rust VGA Kernel - x86_64 / modo longo
  Boot: Multiboot2 (BIOS)
  Video: VGA texto 80x25 @ 0xB8000
  Memoria: paginas de 4 KiB + demand paging
  Interrupcoes: IDT + PIC + timer IRQ0 (~100 Hz)
==============================================================
Multiboot2: magic=0x36d76289 info=0x00112040

Digite 'help' para ver os comandos.

rinux>
\end{lstlisting}

The value of \code{info} varies on each execution, because it depends on the
allocation made by GRUB. The \code{magic} must always be
\code{0x36d76289}.

\subsection{\code{help} command via PS/2}

\begin{lstlisting}[style=shell,caption={Output of \code{help} after PS/2 input.},label={lst:testhelp}]
rinux> help
Comandos disponiveis:
  help          mostra esta ajuda
  clear         limpa o console ativo
  echo <texto>  imprime o texto
  info          informacoes do kernel
  tty [0..7]    mostra ou troca o console virtual
  getty [0..7]  roda o login em ring 3 no console
  demand        testa demand paging no kernel
  usertest      roda um programa em ring 3 com page fault
  userprot      testa a protecao usuario/kernel
  mem           estatisticas de frames e paginas
  reboot        reinicia o PC
  halt          para a CPU
rinux>
\end{lstlisting}

\subsection{\code{echo} command via PS/2}

\begin{lstlisting}[style=shell,caption={Output of \code{echo ola} after PS/2 input.},label={lst:testecho}]
rinux> echo ola
ola
rinux>
\end{lstlisting}

\subsection{\code{clear} via ANSI}

\begin{lstlisting}[style=shell,caption={Escape sequence observed on COM1.},label={lst:testclear}]
rinux> clear
^[[2J^[[H
rinux>
\end{lstlisting}

The caret notation is how a serial monitor displays the byte \code{0x1B}. On
the VGA console the command blanks all 2,000 cells and returns the cursor to
\code{(0,0)}, with no escape characters printed as text.

\subsection{Serial input}

\begin{lstlisting}[style=shell,caption={Input and output through COM1.},label={lst:testserial}]
echo serial ok
serial ok
rinux>
\end{lstlisting}

The text typed on the serial socket appears on the VGA because the shell echoes
each character and the console writes both to the VGA and to the serial port.

\subsection{Signals and process lifecycle}

A worker created with \code{spawn sig1 1 signal} installs a \code{SIGUSR1}
handler and sleeps. Sending \code{kill 2 USR1} wakes it and the handler runs
in the target's context:

\begin{lstlisting}[style=shell,caption={Custom signal handler.},label={lst:testsignal}]
rinux> spawn sig1 1 signal
tarefa criada: pid=2 nome=sig1 nice=0
[task 2] esperando SIGUSR1 (kill 2 USR1)
rinux> kill 2 USR1
pid=2 sinal=10 (SIGUSR1) enviado
[task 2] handler: recebi SIGUSR1 (10)
\end{lstlisting}

Default actions, stopping and reaping are exercised with:

\begin{lstlisting}[style=shell,caption={Default actions, SIGSTOP/SIGCONT and wait.},label={lst:testlifecycle}]
rinux> spawn stop1 1 sleep
tarefa criada: pid=2 nome=stop1 nice=0
rinux> kill 2 STOP
pid=2 sinal=19 (SIGSTOP) enviado
rinux> ps
*    1     0  rodando        ...  shell
     2     1  parada         ...  stop1
rinux> kill 2 CONT
rinux> kill 2 TERM
pid=2 sinal=15 (SIGTERM) enviado
rinux> wait 2
pid=2 terminado por SIGTERM (15), status=143
rinux> spawn short 1 exit
tarefa criada: pid=3 nome=short nice=0
[task 3] encerrando com status 42 em 0,5 s
rinux> wait 3
pid=3 status=42
\end{lstlisting}

Blocking is symmetric: \code{sigblock <pid> USR1 on} keeps \code{SIGUSR1}
pending (visible in \code{sig <pid>} and \code{/proc/signals}) until
\code{sigblock <pid> USR1 off}; the next safe point then runs the handler.
\code{kill <pid> 0} reports that a zombie still exists, and after
\code{wait <pid>} the same call fails with \emph{tarefa nao encontrada}.

\section{Headless test script}

The procedure below can be adapted for CI. It starts QEMU without a
graphical interface, sends a command over the serial port and reads the
output.

\begin{lstlisting}[language=Python,caption={Headless test skeleton with a serial socket.},label={lst:testscript}]
import socket, subprocess, time, os

ISO = "vga-kernel.iso"
SER = "/tmp/vga-serial.sock"
MON = "/tmp/vga-mon.sock"

for path in (SER, MON):
    try: os.unlink(path)
    except FileNotFoundError: pass

qemu = subprocess.Popen([
    "qemu-system-x86_64", "-cdrom", ISO, "-boot", "order=d",
    "-m", "128M", "-display", "none",
    "-serial", f"unix:{SER},server=on,wait=off",
    "-monitor", f"unix:{MON},server=on,wait=off",
])

time.sleep(7)
ser = socket.socket(socket.AF_UNIX, socket.SOCK_STREAM)
ser.connect(SER)
ser.sendall(b"echo serial ok\r")
time.sleep(1)
data = ser.recv(65536).decode(errors="replace")
assert "serial ok" in data
ser.close()
qemu.terminate()
qemu.wait(timeout=10)
\end{lstlisting}

To check the VGA, connect to the monitor, send
\code{xp /4000bx 0xb8000}, interpret the pairs (character, attribute) and
reconstruct the lines. The PS/2 test uses \code{sendkey}:

\begin{lstlisting}[style=shell,caption={Sending keys through the QEMU monitor.},label={lst:sendkey}]
sendkey h
sendkey e
sendkey l
sendkey p
sendkey ret
\end{lstlisting}

\begin{atencao}[Do not send \code{quit} before finishing]
The \code{quit} command in the monitor terminates QEMU. If the script still
needs to inspect the VGA or send more keys, use a separate connection for each
command and only terminate the process at the end.
\end{atencao}

\section{Build and format tests}

\begin{lstlisting}[style=shell,caption={Structural validation commands.},label={lst:testbuild}]
make clean && make
grub-file --is-x86-multiboot2 target/x86_64-unknown-none/release/vga-kernel
readelf -h target/x86_64-unknown-none/release/vga-kernel
readelf -l target/x86_64-unknown-none/release/vga-kernel
xorriso -indev vga-kernel.iso -report_el_torito plain
\end{lstlisting}

Expected results:

\begin{itemize}
  \item \code{grub-file} returns code 0;
  \item \code{readelf -h} shows \code{Type: EXEC}, \code{Machine: Advanced
        Micro Devices X86-64} and \code{Entry point address: 0x100018};
  \item \code{readelf -l} shows a low segment at \code{0x100000} and
        high segments at \code{0xFFFFFFFF8010...};
  \item \code{xorriso} reports \emph{El Torito} with
        \code{boot\_hybrid.img}.
\end{itemize}

\section{Test limitations}

\begin{itemize}
  \item There is no automated coverage of all keyboard paths
        (Shift, Tab, Backspace, extended keys).
  \item The PS/2 test depends on QEMU's \code{sendkey} and on the monitor's
        key layout; not all characters are easy to simulate.
  \item Reading the VGA through the monitor depends on the dump format and on
        the \code{0xB8000} address.
  \item Behavior on physical hardware may differ from QEMU,
        especially on UEFI without CSM, USB keyboards and chipsets without 8042.
  \item \code{reboot} and \code{halt} were verified manually; there is no
        automated reset test.
  \item There are no stress, timing or concurrency tests.
\end{itemize}

\section{Acceptance checklist}

\begin{itemize}
  \item[$\square$] \code{make} generates the ELF and the ISO without errors.
  \item[$\square$] \code{grub-file} recognizes Multiboot2.
  \item[$\square$] QEMU reaches the banner and the \code{rinux>} prompt.
  \item[$\square$] \code{help}, \code{echo}, \code{info} and \code{clear}
        work through the PS/2 keyboard.
  \item[$\square$] Serial input works.
  \item[$\square$] \code{info} shows the timer and \code{mem} shows
        increasing \emph{ticks}.
  \item[$\square$] \code{demand} maps the kernel page after
        \code{\#PF}.
  \item[$\square$] \code{usertest} raises two page faults and returns to the shell.
  \item[$\square$] \code{userprot} is terminated by the U/S protection.
  \item[$\square$] \code{tty 3} switches to \code{tty3} and \code{tty 0}
        returns to \code{tty0}.
  \item[$\square$] \code{getty} accepts \code{root}/\code{rinux} and returns
        to the shell; a wrong password exits with status 1.
  \item[$\square$] A \code{SIGUSR1} handler installed with \code{spawn ...
        signal} runs after \code{kill <pid> USR1}.
  \item[$\square$] \code{SIGTERM}/\code{SIGKILL} terminate a worker and
        \code{SIGSTOP}/\code{SIGCONT} stop and resume it.
  \item[$\square$] A blocked signal stays pending until \code{sigblock ... off}.
  \item[$\square$] \code{wait} reaps status 42 and signal status 143 and shows
        them; \code{/proc/signals} matches \code{sig}.
  \item[$\square$] \code{halt} stops the CPU and \code{reboot} restarts (when the
        environment supports it).
  \item[$\square$] The LaTeX documentation compiles and generates the PDF.
\end{itemize}
